\begin{proofbox}
	\textbf{\textcolor{CProof}{思路提示}}
	利用特殊三棱锥（从长方体一角切出）验证体积公式，再借助祖暅原理（等底等高锥体体积相等）推广到任意锥体。

	\medskip
	\textbf{\textcolor{CProof}{正式推导}}
	\begin{description}[style=nextline, leftmargin=2.8em, labelsep=0.5em, itemsep=0.55em]
		\item[\textbf{步骤一（特殊三棱锥）：}]
			构造一个三条棱两两垂直的三棱锥$P-ABC$，其中$PA=a$，$PB=b$，$PC=c$。取点$C$为顶点，底面为$\triangle PAB$，则$PC\perp$平面$PAB$，故高$h=PC=c$，底面积$S_{\triangle PAB}=\frac{1}{2}ab$。该三棱锥可看作从长方体中切出的一个角。已知长方体体积为$abc$，且长方体可分割为$6$个体积相等的三棱锥，故$V=\frac{1}{6}abc$。代入得：$V=\frac{1}{3}\cdot\frac{1}{2}ab\cdot c=\frac{1}{3}Sh$，公式成立。
			\par\textbf{理由：}长方体分割可得三棱锥体积，验证了公式在特殊情形下成立。
		\item[\textbf{步骤二（等积推广）：}]
			由祖暅原理，等底等高的两个锥体体积相等。因此，任意一个三棱锥，若其底面积也为$S$、高也为$h$，则其体积必等于上述特殊三棱锥的体积，即$\frac{1}{3}Sh$。故所有三棱锥的体积公式均为$V=\frac{1}{3}Sh$。
			\par\textbf{理由：}祖暅原理保证体积只与底面积和高有关，与形状无关。
		\item[\textbf{步骤三（任意锥体）：}]
			任意$n$棱锥可以被通过顶点的分割线分成$(n-2)$个三棱锥，每个小三棱锥的高与原棱锥的高相同，底面积之和等于原底面积$S$。因此原棱锥体积$V=\frac{1}{3}Sh$。圆锥可以看作底面正多边形边数无限增加时的极限，或用祖暅原理与等底等高三棱锥比较，也服从该公式。
			\par\textbf{理由：}分割法及极限思想（或祖暅原理）确保统一公式。
	\end{description}

	\medskip
	\textbf{\textcolor{CProof}{结论回扣}}
	综上，无论棱锥还是圆锥，只要底面积为$S$、高为$h$，其体积均为$V=\frac{1}{3}Sh$。
\end{proofbox}
